\chapter*{阅读路径}
\addcontentsline{toc}{chapter}{阅读路径}

三部分可以相对独立阅读。代数与分析部分遵循“群与表示 \(\to\) 结合代数与 Clifford/Spin \(\to\) Hilbert 空间与广义本征结构 \(\to\) 自伴算子、谱测度与 resolvent \(\to\) 分布/Fourier \(\to\) Fredholm 与逆问题 \(\to\) 非线性演化与可积结构”。其中需要 Sturm--Liouville 或 Green 函数时，先使用一般闭算子和谱演算，再下降到微分算子模型。

微分几何部分遵循“切空间与流 \(\to\) Lie 导数 \(\to\) 标架/余标架 \(\to\) 向量丛联络 \(\to\) 挠率/曲率 \(\to\) 度量/Levi--Civita \(\to\) 指数映射与 Jacobi 场 \(\to\) Killing/Hodge \(\to\) Frobenius 与 holonomy \(\to\) 几何演化 \(\to\) moving frame/Gauss--Codazzi--Ricci”。统计部分遵循“概率模型 \(\to\) LLN/CLT 与抽样分布 \(\to\) sufficiency/completeness \(\to\) likelihood/Fisher information 与 Bayes \(\to\) 检验理论 \(\to\) 一般线性模型”。

若只为物理课程补充数学工具，可以按依赖关系跳读，但不建议跨过定义直接使用后面的坐标公式：例如使用旋量前至少读 Clifford/Spin 的普适性质；使用曲率分量前至少读一般联络与 Cartan 方程；使用 \(F\) 检验前至少读投影、充分统计量与正态样本分解。
